% Provenance (never printed): see sources/outline_sources.md and sources/book1/BRIEFS.md (g7/06).
\chapter{Project: A Quiz Game}\label{ch:g7:project-quiz-game}

This chapter builds one \omterm{def:g4:programs-and-operating-system:program}{program} from start to finish, a quiz game, in five
versions. Each version runs, each is tested, and each adds one thing to the
one before. Here is what the finished game must do:

\begin{quote}
\itshape A quiz asks ten questions, read from \omterm{def:g7:lists:list}{lists}. The player may type an
answer in capitals or small letters, with extra spaces around it. A right
answer gives one point; after a wrong one the player has a second chance,
worth half a point. The score is shown at the end. The quiz has an easy level
and a hard level.
\end{quote}

\begin{recall}
A \omterm{def:g7:functions:function}{function} takes \omterm{def:g7:functions:parameter}{parameters} and returns a \omterm{def:g6:variables-and-input:value}{value} (\cref{ch:g7:functions}). A
\omterm{def:g7:lists:list}{list} keeps \omterm{def:g6:variables-and-input:value}{values} in order, with indices from 0 (\cref{ch:g7:lists}).
\code{strip} and \code{lower} return a cleaned copy of a \omterm{def:g7:strings:string}{string}
(\cref{ch:g7:strings}). A \code{while} \omterm{def:g4:algorithms-repeat-and-decide:loop}{loop} repeats until its \omterm{def:g4:algorithms-repeat-and-decide:condition}{condition} is
false (\cref{ch:g6:while-loops}). \omterm{def:g7:finding-and-fixing-bugs:bug}{Bugs} are found by reproducing, shrinking and
changing one thing at a time (\cref{ch:g7:finding-and-fixing-bugs}).
\end{recall}

\section{What must it do?}

The paragraph above is written for people. Before writing code, turn it into
a numbered \omterm{def:g7:lists:list}{list} of things to check, each one small enough to test:

\begin{enumerate}
  \item The questions and their answers are kept in two \omterm{def:g7:lists:list}{lists} of the same
    length.
  \item An answer is accepted whatever its capitals and its spaces at the
    ends.
  \item A right first answer gives 1 point, a right second answer gives
    0.5, a wrong second answer gives 0.
  \item Every question is asked once, in order, and the score is printed at
    the end.
  \item The player chooses the easy or the hard level, and is asked again if
    the level typed is neither.
  \item The \omterm{def:g4:programs-and-operating-system:program}{program} ends with one line in a fixed format, which another
    \omterm{def:g4:programs-and-operating-system:program}{program} could read.
\end{enumerate}

Then plan the \omterm{def:g7:functions:function}{functions}. Each one does one job and can be tested alone:
\code{normalise} cleans a typed answer, \code{ask} asks one question,
\code{play} asks them all, \code{choose_level} asks for the level, and
\code{main} puts them together.

\begin{omfigure}
\begin{tikzpicture}[font=\footnotesize\ttfamily, f/.style={draw=black!60, rounded corners, fill=ocBlue!10, minimum width=2.4cm, minimum height=0.65cm}]
  \node[f, fill=ocYellow!30] (m) at (0,0) {main};
  \node[f] (c) at (-3.4,-1.4) {choose\_level};
  \node[f] (p) at (0,-1.4) {play};
  \node[f] (r) at (3.4,-1.4) {result\_line};
  \node[f] (a) at (0,-2.8) {ask};
  \node[f] (n) at (-3.4,-2.8) {normalise};
  \node[f, fill=white] (i) at (3.4,-2.8) {input};
  \draw[-{Stealth[length=2mm]}, thick] (m) -- (c);
  \draw[-{Stealth[length=2mm]}, thick] (m) -- (p);
  \draw[-{Stealth[length=2mm]}, thick] (m) -- (r);
  \draw[-{Stealth[length=2mm]}, thick] (p) -- (a);
  \draw[-{Stealth[length=2mm]}, thick] (a) -- (n);
  \draw[-{Stealth[length=2mm]}, thick] (c) -- (n);
  \draw[-{Stealth[length=2mm]}, thick] (a) -- (i);
\end{tikzpicture}

{\small The plan of the finished \omterm{def:g4:programs-and-operating-system:program}{program}: an arrow goes from a \omterm{def:g7:functions:function}{function} to
each \omterm{def:g7:functions:function}{function} it calls. \code{input} is Python's own.}
\end{omfigure}

\section{Questions in lists}

The first version only holds the questions and the answers, in two \omterm{def:g7:lists:list}{lists} that
go together: \code{answers[i]} is the answer to \code{questions[i]}. Answers
are written in small letters, so that they can be compared with what the
player types, once cleaned.

\omcode{code/grade-7/06-project-quiz-game/listings/quiz_v1.py}{1}{19}{Version 1:
the questions and their answers.}\label{lst:g7:project-quiz-game:v1}

\Needspace*{11\baselineskip}
\omoutput[1-5,10-11]{code/grade-7/06-project-quiz-game/outputs/quiz_v1.txt}{The
first four pairs and the last one.}\label{out:g7:project-quiz-game:v1}

\begin{omfigure}
\begin{tikzpicture}[x=1cm, y=1cm, font=\scriptsize\ttfamily]
  \foreach \q/\a [count=\i from 0] in {France/paris, spider/8, 7 x 8/56, largest/jupiter, bees/honey} {
    \draw[black!60, fill=ocBlue!10] (2.1*\i,0.8) rectangle ++(2.1,0.6);
    \node at (2.1*\i+1.05,1.1) {\q?};
    \draw[black!60, fill=ocYellow!20] (2.1*\i,0) rectangle ++(2.1,0.6);
    \node at (2.1*\i+1.05,0.3) {\a};
    \node[ocRed!80!black, font=\scriptsize\sffamily] at (2.1*\i+1.05,-0.25) {\i};
  }
  \node[anchor=east, font=\footnotesize\sffamily] at (-0.1,1.1) {questions};
  \node[anchor=east, font=\footnotesize\sffamily] at (-0.1,0.3) {answers};
  \node[font=\footnotesize\sffamily] at (10.9,0.7) {\dots};
\end{tikzpicture}

{\small Two parallel \omterm{def:g7:lists:list}{lists}: the same \omterm{def:g7:lists:list}{index} picks a question and its answer
(questions shortened here).}
\end{omfigure}

The check at line 16 catches a \omterm{def:g7:lists:list}{list} that is one \omterm{def:g7:lists:list}{item} short, before the game
goes wrong. From now on, the \omterm{def:g7:lists:list}{lists} live in a \omterm{def:g3:files-and-folders:file}{file} of their own,
\code{questions.py}, with a second pair of \omterm{def:g7:lists:list}{lists} for the hard level; the next
versions import them.

\section{Asking one question}

The player may type \code{Paris}, \code{PARIS} or \code{" paris "}.
\code{normalise} removes the spaces at the ends and puts the letters in small
type, so that all of these become \code{"paris"}.

\omcode{code/grade-7/06-project-quiz-game/listings/quiz_v2.py}{1}{18}{Version 2:
asking one question.}\label{lst:g7:project-quiz-game:v2}

\Needspace*{7\baselineskip}
\omoutput{code/grade-7/06-project-quiz-game/outputs/quiz_v2.txt}{An answer in
capitals, with spaces, is accepted.}\label{out:g7:project-quiz-game:v2}

A test of the chapter calls \code{ask} without a keyboard: it replaces
\code{input} with a \omterm{def:g7:functions:function}{function} that returns scripted answers, then checks the
points that \code{ask} returns.

\section{The whole quiz}

\code{play} walks through the two \omterm{def:g7:lists:list}{lists} with one \omterm{def:g7:lists:list}{index} and adds up the
points.

\omcode{code/grade-7/06-project-quiz-game/listings/quiz_v3.py}{1}{17}{Version 3:
the whole quiz.}\label{lst:g7:project-quiz-game:v3}

\Needspace*{13\baselineskip}
\omoutput[1-5,20-23]{code/grade-7/06-project-quiz-game/outputs/quiz_v3.txt}{A
perfect game.}\label{out:g7:project-quiz-game:v3}

\code{quiz_v3_poor.py} imports \code{play} and plays the same quiz with poorer
answers:

\Needspace*{12\baselineskip}
\omoutput[1-5,20-22]{code/grade-7/06-project-quiz-game/outputs/quiz_v3_poor.txt}{A
poorer game: 3 out of 10.}\label{out:g7:project-quiz-game:v3-poor}

\section{Second chances and levels}

A wrong first answer now gets a second try. \code{ask} returns 1, 0.5 or 0,
so the score starts at \code{0.0} and is a decimal number. A second \omterm{def:g7:functions:function}{function}
asks for the level until the player types one of the two.

\omcode{code/grade-7/06-project-quiz-game/listings/quiz_v4.py}{1}{40}{Version 4:
second chances and two levels.}\label{lst:g7:project-quiz-game:v4}

\omoutput[1-10,24-27]{code/grade-7/06-project-quiz-game/outputs/quiz_v4.txt}{A
hard game with a mistyped level and three second
chances.}\label{out:g7:project-quiz-game:v4}

\begin{omfigure}
\begin{tikzpicture}[font=\footnotesize, b/.style={draw=black!60, rounded corners, fill=ocBlue!10, align=center, minimum height=0.7cm}, d/.style={draw=black!60, diamond, aspect=2.2, fill=ocYellow!25, inner sep=1pt, align=center}]
  \node[b] (q) at (0,0) {ask the question};
  \node[d] (r1) at (0,-1.4) {right?};
  \node[b, fill=ocMint!30] (p1) at (3.6,-1.4) {1 point};
  \node[b] (q2) at (0,-2.8) {ask again};
  \node[d] (r2) at (0,-4.2) {right?};
  \node[b, fill=ocMint!30] (p2) at (3.6,-4.2) {0.5 point};
  \node[b, fill=ocRed!10] (p0) at (-3.6,-4.2) {0 points,\\ show the answer};
  \draw[-{Stealth[length=2mm]}, thick] (q) -- (r1);
  \draw[-{Stealth[length=2mm]}, thick] (r1) -- node[above] {yes} (p1);
  \draw[-{Stealth[length=2mm]}, thick] (r1) -- node[left] {no} (q2);
  \draw[-{Stealth[length=2mm]}, thick] (q2) -- (r2);
  \draw[-{Stealth[length=2mm]}, thick] (r2) -- node[above] {yes} (p2);
  \draw[-{Stealth[length=2mm]}, thick] (r2) -- node[above] {no} (p0);
\end{tikzpicture}

{\small One turn of the game, as \code{ask} plays it in version 4.}
\end{omfigure}

\section{Improving it}

The last version adds the final line of the \omterm{def:g7:lists:list}{list} of requirements: a result
in a fixed format, which a person can read and a \omterm{def:g4:programs-and-operating-system:program}{program} can split into its
parts. Shuffling the order of the questions would make the game better, but
it needs random numbers, which come in grade 9.

\omcode{code/grade-7/06-project-quiz-game/listings/quiz_v5.py}{1}{23}{Version 5:
the finished game.}\label{lst:g7:project-quiz-game:v5}

\omoutput[1-7,17-26]{code/grade-7/06-project-quiz-game/outputs/quiz_v5.txt}{An easy
game, scored 8.0 out of 10.}\label{out:g7:project-quiz-game:v5}

\begin{omfigure}
\begin{tikzpicture}[x=1cm, y=1cm, font=\footnotesize]
  \draw[-{Stealth[length=2mm]}, thick] (0,0) -- (12.3,0);
  \foreach \v/\t/\x in {v1/{lists, lengths}/0.8, v2/{normalise, ask}/3.2, v3/{play, two games}/5.6, v4/{half points, levels}/8.0, v5/{result line}/10.4} {
    \fill[ocBlue] (\x,0) circle (0.09);
    \node[above, font=\footnotesize\ttfamily] at (\x,0.15) {\v};
    \node[below, align=center, text width=2.2cm] at (\x,-0.15) {\t};
  }
\end{tikzpicture}

{\small The five versions, with what each one adds.}
\end{omfigure}

\section{Exercises}

\begin{exercise}[$\star$]\label{exo:g7:project-quiz-game:1}
What does \code{normalise("  Paris ")} return? And
\code{normalise("NEW YORK")}?
\end{exercise}

\begin{exercise}[$\star$]\label{exo:g7:project-quiz-game:2}
In version 3, a player types \code{Paris}, \code{eight}, \code{56},
\code{jupiter}, \code{" Honey "}, \code{60}, \code{hot}, \code{6},
\code{green} and \code{63}. What score is printed?
\end{exercise}

\begin{exercise}[$\star$]\label{exo:g7:project-quiz-game:3}
Change \code{ask} of version 4 so that typing \code{skip} moves on at once,
with 0 points and no second chance.
\end{exercise}

\begin{exercise}[$\star$]\label{exo:g7:project-quiz-game:4}
After an edit, the quiz says that every answer from question 5 on is wrong,
and stops with an \code{IndexError} at question 10. Someone deleted one \omterm{def:g7:lists:list}{item}
from a \omterm{def:g7:lists:list}{list}. Which \omterm{def:g7:lists:list}{list}, and which line of version 1 would have warned about
it?
\end{exercise}

\begin{exercise}[$\star\star$]\label{exo:g7:project-quiz-game:5}
Make version 5 print the score as a percentage as well.
\end{exercise}

\begin{exercise}[$\star\star$]\label{exo:g7:project-quiz-game:6}
This \code{ask} gives 1.5 points for \code{paris} typed twice. Why? Fix it.

\omcode{code/grade-7/06-project-quiz-game/exercises/ask_bug.py}{5}{11}{Too
many points.}\label{lst:g7:project-quiz-game:ask-bug}
\end{exercise}

\begin{exercise}[$\star\star$]\label{exo:g7:project-quiz-game:7}
Add a third level, \code{medium}. Which parts of the \omterm{def:g4:programs-and-operating-system:program}{program} change?
\end{exercise}

\begin{exercise}[$\star\star$]\label{exo:g7:project-quiz-game:8}
How many lines does a game of version 5 show on the screen at most, when the
player types a valid level at once? Why is there no maximum otherwise?
\end{exercise}

\begin{exercise}[$\star\star$]\label{exo:g7:project-quiz-game:9}
Write a \code{play} that stops early when the player types \code{quit}, and
returns the score so far. Use a \code{while} \omterm{def:g4:algorithms-repeat-and-decide:loop}{loop}.
\end{exercise}

\begin{exercise}[$\star\star\star$]\label{exo:g7:project-quiz-game:10}
Write a \omterm{def:g7:functions:function}{function} \code{longest_streak(points)} that takes the points of each
question, such as \code{[1, 0.5, 1, 1, 1, 0, 1, 1, 0.5, 1]}, and returns the
longest run of right first answers.
\end{exercise}

\begin{exercise}[$\star\star\star$]\label{exo:g7:project-quiz-game:11}
What are the highest and the lowest possible scores in a game of ten
questions? Can every score in between, in steps of 0.5, be reached? How many
different scores are there?
\end{exercise}

\begin{exercise}[$\star\star\star$]\label{exo:g7:project-quiz-game:12}
A spider has \code{8} legs, but a player who types \code{eight} is right too.
Change the answers into a \omterm{def:g7:lists:list}{list} of \omterm{def:g7:lists:list}{lists}, with every accepted answer, and
change \code{ask} to match.
\end{exercise}
